% ============================================================================
%  Preámbulo
% ============================================================================

% ---- Idioma y tipografía ---------------------------------------------------
\usepackage[T1]{fontenc}
\usepackage[utf8]{inputenc}
\usepackage[spanish,es-noshorthands,es-tabla,es-nodecimaldot]{babel}
\usepackage{lmodern}
\usepackage{microtype}
\usepackage{csquotes}

% ---- Página ----------------------------------------------------------------
\usepackage[margin=2.6cm,headheight=14pt]{geometry}
\usepackage{fancyhdr}
\pagestyle{fancy}
\fancyhf{}
\fancyhead[L]{\small\nouppercase{\leftmark}}
\fancyhead[R]{\small Algoritmos del núcleo de Tukit}
\fancyfoot[C]{\small\thepage}
\renewcommand{\headrulewidth}{0.4pt}
\setlength{\parskip}{0.35em}

% ---- Matemáticas -----------------------------------------------------------
\usepackage{amsmath,amssymb,amsthm,mathtools}
\usepackage{siunitx}
\sisetup{output-decimal-marker={,}, group-separator={.}, group-minimum-digits=4}

% ---- Tablas, figuras, listas ----------------------------------------------
\usepackage{booktabs,multirow,threeparttable,array}
\newcolumntype{L}[1]{>{\raggedright\arraybackslash}p{#1}}
\usepackage{graphicx}
\usepackage{float}
\usepackage[font=small,labelfont=bf]{caption}
\usepackage[inline]{enumitem}
\setlist{itemsep=0.15em,topsep=0.3em}

% ---- Pseudocódigo ----------------------------------------------------------
\usepackage{algorithm}
\usepackage[noend]{algpseudocode}
\floatname{algorithm}{Algoritmo}
\renewcommand{\listalgorithmname}{Lista de algoritmos}
\algrenewcommand\algorithmicrequire{\textbf{Entrada:}}
\algrenewcommand\algorithmicensure{\textbf{Salida:}}
\algrenewcommand\algorithmicif{\textbf{si}}
\algrenewcommand\algorithmicthen{\textbf{entonces}}
\algrenewcommand\algorithmicelse{\textbf{si no}}
\algrenewcommand\algorithmicfor{\textbf{para}}
\algrenewcommand\algorithmicforall{\textbf{para todo}}
\algrenewcommand\algorithmicdo{\textbf{hacer}}
\algrenewcommand\algorithmicwhile{\textbf{mientras}}
\algrenewcommand\algorithmicreturn{\textbf{devolver}}
\algrenewcommand\algorithmicfunction{\textbf{función}}
\algrenewcommand\algorithmicprocedure{\textbf{procedimiento}}
\algrenewcommand\algorithmicrepeat{\textbf{repetir}}
\algrenewcommand\algorithmicuntil{\textbf{hasta}}
\algnewcommand\algorithmicand{\textbf{y}}
\algnewcommand\algorithmicor{\textbf{o}}
\algnewcommand\AND{\algorithmicand}
\algnewcommand\OR{\algorithmicor}
\algnewcommand{\LineComment}[1]{\State \(\triangleright\) \textit{#1}}
\algrenewcommand\alglinenumber[1]{\scriptsize #1:}

% ---- Diagramas y gráficas -------------------------------------------------
\usepackage{tikz}
\usetikzlibrary{arrows.meta,positioning,shapes.geometric,fit,calc,backgrounds}
\usepackage{pgfplots}
\pgfplotsset{compat=1.18}
\usepgfplotslibrary{groupplots}

% ---- Colores y bloques (en la línea de las notas de clase) ----------------
\usepackage{xcolor}
\definecolor{azul}{RGB}{40,40,150}
\definecolor{azulclaro}{RGB}{232,232,246}
\definecolor{verde}{RGB}{20,110,40}
\definecolor{verdeclaro}{RGB}{230,243,232}
\definecolor{rojo}{RGB}{170,40,40}
\definecolor{rojoclaro}{RGB}{248,232,232}
\definecolor{gris}{RGB}{90,90,90}
\definecolor{grisclaro}{RGB}{242,242,242}

\usepackage[most]{tcolorbox}
\tcbset{
  bloque/.style={enhanced, breakable, arc=1.5mm, boxrule=0.6pt,
    left=2mm, right=2mm, top=1mm, bottom=1mm,
    fonttitle=\bfseries\small, coltitle=white},
}

% Caja para registrar la evolución de un algoritmo: antes / defecto / cambio.
\newtcolorbox{evolucion}[1][]{bloque, colback=grisclaro, colframe=gris,
  title={Evolución#1}}
% Caja de advertencia o de hallazgo crítico (el título es obligatorio).
\newtcolorbox{hallazgo}[1]{bloque, colback=rojoclaro, colframe=rojo,
  title={#1}}

% ---- Referencias cruzadas --------------------------------------------------
\usepackage{xurl}
\usepackage[hidelinks,pdfusetitle,bookmarksnumbered]{hyperref}
\usepackage[spanish,nameinlink,noabbrev]{cleveref}

% Teoremas con numeración por sección. Se definen DESPUÉS de cleveref para que
% \cref distinga entornos que comparten contador (teorema, proposición, ...).
% Cada entorno tiene un contador alias del de `definicion`: comparten la
% numeración, pero cleveref sabe cómo llamar a cada uno.
\usepackage{aliascnt}
\newcommand{\nuevoenunciado}[2]{%
  \newaliascnt{#1}{definicion}%
  \newtheorem{#1}[#1]{#2}%
  \aliascntresetthe{#1}}
\theoremstyle{definition}
\newtheorem{definicion}{Definición}[section]
\nuevoenunciado{ejemplo}{Ejemplo}
\nuevoenunciado{observacion}{Observación}
\theoremstyle{plain}
\nuevoenunciado{proposicion}{Proposición}
\nuevoenunciado{lema}{Lema}
\nuevoenunciado{teorema}{Teorema}
\nuevoenunciado{corolario}{Corolario}
\renewcommand{\qedsymbol}{$\blacksquare$}

% Cajas de color alrededor de los enunciados principales.
\tcolorboxenvironment{definicion}{bloque, colback=azulclaro, colframe=azul}
\tcolorboxenvironment{proposicion}{bloque, colback=verdeclaro, colframe=verde}
\tcolorboxenvironment{teorema}{bloque, colback=verdeclaro, colframe=verde}
\tcolorboxenvironment{lema}{bloque, colback=verdeclaro, colframe=verde}
\crefname{section}{sección}{secciones}
\crefname{subsection}{sección}{secciones}
\crefname{subsubsection}{sección}{secciones}
\crefname{appendix}{apéndice}{apéndices}
\crefname{table}{tabla}{tablas}
\crefname{figure}{figura}{figuras}
\crefname{equation}{ecuación}{ecuaciones}
\crefname{enumi}{requisito}{requisitos}
\crefname{definicion}{definición}{definiciones}
\crefname{proposicion}{proposición}{proposiciones}
\crefname{lema}{lema}{lemas}
\crefname{teorema}{teorema}{teoremas}
\crefname{corolario}{corolario}{corolarios}
\crefname{observacion}{observación}{observaciones}
\crefname{ejemplo}{ejemplo}{ejemplos}
\crefname{algorithm}{algoritmo}{algoritmos}
\crefname{ALG@line}{línea}{líneas}

% ---- Notación ----------------------------------------------------------------
\newcommand{\N}{\mathbb{N}}
\newcommand{\R}{\mathbb{R}}
\newcommand{\Z}{\mathbb{Z}}
\newcommand{\Cat}{\mathcal{C}}          % conceptos
\newcommand{\Prod}{\mathcal{P}}         % productos
\newcommand{\Tipos}{\mathcal{T}}        % tipos de renglón
\newcommand{\Est}{\mathcal{E}}          % estados de resolución
\newcommand{\Dec}{\mathcal{D}}          % líneas decidibles
\newcommand{\Esp}{\Omega}               % espacio de búsqueda
\newcommand{\Front}{\mathcal{F}}        % frente
\newcommand{\nivel}{\nu}
\newcommand{\precio}{\pi}
\newcommand{\emp}{\upsilon}             % unidades por empaque
\newcommand{\techo}[1]{\left\lceil #1 \right\rceil}
\newcommand{\piso}[1]{\left\lfloor #1 \right\rfloor}
\newcommand{\abs}[1]{\left\lvert #1 \right\rvert}
\newcommand{\card}[1]{\left\lvert #1 \right\rvert}
\newcommand{\ind}[1]{\mathbf{1}\!\left[#1\right]}
\newcommand{\domina}{\prec}
\newcommand{\sinprecio}{\bot}
\newcommand{\cod}[1]{\texttt{#1}}
\DeclareMathOperator*{\argmin}{arg\,min}
\DeclareMathOperator*{\argmax}{arg\,max}
\DeclareMathOperator*{\lexmin}{lex\,min}
\DeclareMathOperator*{\lexmax}{lex\,max}
\DeclareMathOperator{\sing}{sing}
\DeclareMathOperator{\cab}{cab}
\DeclareMathOperator{\sig}{sig}
